\documentclass[11pt]{article}
\usepackage{amsmath,amssymb,amsthm,enumerate,nicefrac,fancyhdr,hyperref,graphicx,adjustbox}
\hypersetup{colorlinks=true,urlcolor=blue,citecolor=blue,linkcolor=blue}
\usepackage[left=2.6cm, right=2.6cm, top=1.5cm, includehead, includefoot]{geometry}
\usepackage[dvipsnames]{xcolor}
\usepackage[d]{esvect}

%% commands
%% useful macros [add to them as needed]
% sets
\newcommand{\C}{{\mathbb{C}}} 
\newcommand{\N}{{\mathbb{N}}}
\newcommand{\Q}{{\mathbb{Q}}}
\newcommand{\R}{{\mathbb{R}}}
\newcommand{\Z}{{\mathbb{Z}}}
\newcommand{\F}{{\mathbb{F}}}

% bases
\newcommand{\mA}{\mathcal{A}}
\newcommand{\mB}{\mathcal{B}}
\newcommand{\mC}{\mathcal{C}}
\newcommand{\mD}{\mathcal{D}}
\newcommand{\mE}{\mathcal{E}}
\newcommand{\mL}{\mathcal{L}}
\newcommand{\mM}{\mathcal{M}}
\newcommand{\mO}{\mathcal{O}}
\newcommand{\mP}{\mathcal{P}}
\newcommand{\mS}{\mathcal{S}}
\newcommand{\mT}{\mathcal{T}}

% linear algebra
\newcommand{\diag}{\operatorname{diag}}
\newcommand{\adj}{\operatorname{adj}}
\newcommand{\rank}{\operatorname{rank}}
\newcommand{\spn}{\operatorname{Span}}
\newcommand{\proj}{\operatorname{proj}}
\newcommand{\prp}{\operatorname{perp}}
\newcommand{\refl}{\operatorname{refl}}
\newcommand{\tr}{\operatorname{tr}}
\newcommand{\nul}{\operatorname{Null}}
\newcommand{\nully}{\operatorname{nullity}}
\newcommand{\range}{\operatorname{Range}}
\renewcommand{\ker}{\operatorname{Ker}}
\newcommand{\col}{\operatorname{Col}}
\newcommand{\row}{\operatorname{Row}}
\newcommand{\cof}{\operatorname{cof}}
\newcommand{\Num}{\operatorname{Num}}
\newcommand{\Id}{\operatorname{Id}}
\newcommand{\ipb}{\langle \thinspace, \rangle}
\newcommand{\ip}[2]{\left\langle #1, #2\right\rangle} % inner products
\newcommand{\M}[2]{M_{#1\times #2}(\F)}
\newcommand{\RREF}{\operatorname{RREF}}
\newcommand{\cv}[1]{\begin{bmatrix} #1 \end{bmatrix}}
\newenvironment{amatrix}[1]{\left[\begin{array}{@{}*{\numexpr#1-1}{c}|c@{}}}{\end{array}\right]}
\newcommand{\am}[2]{\begin{amatrix}{#1} #2 \end{amatrix}}

% vectors
\newcommand{\vzero}{\vv{0}}
\newcommand{\va}{\vv{a}}
\newcommand{\vb}{\vv{b}}
\newcommand{\vc}{\vv{c}}
\newcommand{\vd}{\vv{d}}
\newcommand{\ve}{\vv{e}}
\newcommand{\vf}{\vv{f}}
\newcommand{\vg}{\vv{g}}
\newcommand{\vh}{\vv{h}}
\newcommand{\vl}{\vv{\ell}}
\newcommand{\vm}{\vv{m}}
\newcommand{\vn}{\vv{n}}
\newcommand{\vp}{\vv{p}}
\newcommand{\vq}{\vv{q}}
\newcommand{\vr}{\vv{r}}
\newcommand{\vs}{\vv{s}}
\newcommand{\vt}{\vv{t}}
\newcommand{\vu}{\vv{u}}
\newcommand{\vvv}{{\vv{v}}}
\newcommand{\vw}{\vv{w}}
\newcommand{\vx}{\vv{x}}
\newcommand{\vy}{\vv{y}}
\newcommand{\vz}{\vv{z}}

% display
\newcommand{\ds}{\displaystyle}
\newcommand{\qand}{\quad\text{and}}
\newcommand{\qandq}{\quad\text{and}\quad}
\newcommand{\hint}{\textbf{Hint: }}
\newcommand{\tri}{\triangle}

% misc
\newcommand{\area}{\operatorname{area}}
\newcommand{\vol}{\operatorname{vol}}
\newcommand{\red}[1]{{\color{red} #1}}
\newcommand{\rc}{\red{\checkmark}}

\usepackage{listings}
\usepackage{xcolor}

\definecolor{commentcolor}{rgb}{0.3,0.3,0.3}
\definecolor{keywordcolor}{rgb}{0.1,0.1,0.8}
\definecolor{stringcolor}{rgb}{0.58,0,0.82}

\lstset{
  language=R,
  basicstyle=\ttfamily\footnotesize,
  keywordstyle=\color{blue},
  commentstyle=\color{commentcolor}\itshape,
  stringstyle=\color{stringcolor},
  numbers=none,
  numberstyle=\tiny\color{gray},
  stepnumber=1,
  numbersep=5pt,
  backgroundcolor=\color{white},
  showspaces=false,
  showstringspaces=false,
  showtabs=false,
  frame=none,
  tabsize=2,
  captionpos=b,
  breaklines=true,
  breakatwhitespace=false,
  escapeinside={(*@}{@*)}
}

\title{STAT 341 Notes}
\author{Thomas Liu}
\begin{document}
\maketitle
\tableofcontents

\newpage 
\section{Introduction}
\subsection{The Statistical Narrative}
We direct attention towards populations,
\begin{itemize}
    \item the units which define them 
    \item variates as measures on the units 
\end{itemize}
All genuine populations are finite
\subsection{Population Attributes}
Summaries of finite populations (population attributes)
\begin{itemize}
    \item Define attributes of interest will rely on statistical concepts
    \item Population attributes might involve considerable calculation
    \item No Probability appears for some time as we have the entire population 
    \item Estimate algorithms appear any samples are obtained
\end{itemize}
\subsection{Sample}
All possibe samples 
\begin{itemize}
    \item are explored and characteristics of various population attributes examined within this context
    \item This is done without probability 
\end{itemize}
All possible samples entail a combinatorial explosion for calculation 
\begin{itemize}
    \item the notion of selecting s subset of the samples is introduced
    \item Probability is used for the selection of a sample froom all possible samples 
\end{itemize}
\subsection{Prediction}
\begin{itemize}
    \item Interest lies in predicting the value of a variate (the response variate) given the value of some explanatory variates 
    \item Build a response model that encodes how that prediction is to be carried out 
    \item How should the accuracy be measured?
    \begin{itemize}
        \item Bias-variance tradeoff 
        \item training and test and cross-validation 
    \end{itemize}
\end{itemize}
\newpage 

\section{Populations}
\subsection{Populations}
A population is a finite (though possibly huge) set $\mathcal{P}$ of elements
\begin{itemize}
    \item Elements of a population are called unites $u\in\mathcal{P}$
    \item Variates are functions $x(u), y(u)$, etc. on the which individual unites $u\in\mathcal{P}$. For simplicity we will more often use the notation $x_u$, $y_u$, etc. when referring to the realized values of these variates for the unit $u=1,\cdots,N$ 
\end{itemize}
Population attributes, denoted as $a(\mathcal{P})$. In particular, we will 
\begin{itemize}
    \item Define them 
    \item Consider how to calculate them 
    \item Evaluate some of their (non-sampling) properties, e.g. interpretation of the characteristic being captured, sensitivity to outlying points, etc. 
\end{itemize}
\subsection{Explicitly Defined Population Attributes}
The population is typically a set or collection of units, each with one or more variate that we can measure 
\subsubsection{Population Attributes}
\subsubsection*{Variates}
Variates are characteristics of each unit in the population, and they can take on numerical or categorical values 
\begin{itemize}
    \item The values of variates typically differ from unit to unit 
    \item If we are only interested in the variate $y$'s we might write the population as \[\mathcal{P}=\{y_1,y_2,\cdots,y_N\}\]
\end{itemize}
\subsubsection*{Population Attributes}
Population attributes are summaries describing characteristics of the population 
\begin{itemize}
    \item An attribute is a function applied to the entire population and determined by the variate values observed for each of the population's units \[a(\mathcal{P})=f(y_1,y_2,\cdots,y_N)\]
    \item Some examples of attributes are:
    \begin{itemize}
        \item the population total: \[a(\mathcal{P})=\sum_{u\in\mathcal{P}}y_u\]   
        \item various counts over the population: \[a(\mathcal{P})=\sum_{u\in\mathcal{P}}I_A(y_u)\] where $I_A(y)$ is the indicator function 
        \begin{align*}
            I_A(y) = \begin{cases}
                1 & \text{if } y\in A \\
                0 & \text{otherwise}
            \end{cases}
        \end{align*}
    \end{itemize}
\end{itemize}
\subsubsection*{Location Attributes}
These attributes measure or describe the centre of the distribution of variate values in a dataset
\begin{itemize}
    \item population average/mean: \[a(\mathcal{P})=\overline{y}=\frac{1}{N}\sum_{u\in\mathcal{P}}y_u\]
    \item population proportion: \[a(\mathcal{P})=\frac{1}{N}\sum_{u\in\mathcal{P}}I_A(y_u)\]
    \item Others include mode, median, etc.
\end{itemize}
\subsubsection*{Spread Attributes}
These attributes measure variability or spread of the variate values in a data set
\begin{itemize}
    \item population variance: \[a(\mathcal{P})=\frac{1}{N}\sum_{u\in\mathcal{P}}(y_u-\overline{y})^2\]
    \item population standard deviation: \[a(\mathcal{P})=SD_\mathcal{P}(y)=\sqrt{\frac{1}{N}\sum_{u\in\mathcal{P}}(y_u-\overline{y})^2}\]
    \item coefficient of variation: \[a(\mathcal{P})=\dfrac{SD_\mathcal{P}(y)}{\overline{y}}\]    
    \item Note: population variance or standard deviation could also be defined using $N-1$ in the denominator
    \item Others include range, inter-quartile range, etc.
\end{itemize}
\subsubsection*{Order Statistics}
Population attributes can also be based on an indexed collection of values 
\[y_{(1)}\leq y_{(2)}\leq \cdots \leq y_{(N)}\]
whihch are the variates values $y_u\in\mathcal{P}$ ordered from smallest to largest (including ties)
\subsubsection*{Location Attributes Based on Order Statistics}
\begin{itemize}
    \item population minimum: \[a(\mathcal{P})=\underset{u\in\mathcal{P}}{min}y_u=y_{(1)}\]
    \item population maximum: \[a(\mathcal{P})=\underset{u\in\mathcal{P}}{max}y_u=y_{(N)}\]
    \item population mid range: \[a(\mathcal{P})=\frac{y_{(1)}+y_{(N)}}{2}\]
    \item population median: \[a(\mathcal{P})=\begin{cases}
        \frac{y_{(N/2)}+y_{(N/2+1)}}{2} & \text{if } N \text{ is even} \\
        y_{((N+1)/2)} & \text{if } N \text{ is odd}
    \end{cases}\]
    \item population quartiles:
    \begin{itemize}
        \item $Q_1$ is 25th percentile, or the first quantile 
        \item $Q_2$ is 50th percentile, or the second quantile (median)
        \item $Q_3$ is 75th percentile, or the third quantile
    \end{itemize}
\end{itemize}
\subsubsection*{Variability Attributes Based on Order Statistics}
\begin{itemize}
    \item population range: \[a(\mathcal{P})=y_{(N)}-y_{(1)}\]
    \item population inter-quantile range IQR: \[a(\mathcal{P})=Q_3-Q_1\]
    \item Mdeian Absolute Deviation (MAD) is the median of the absolute differences between eahc $y_u$ and the median \[a(\mathcal{P}) = \underset{u\in\mathcal{P}}{median}|y_u-\underset{u\in\mathcal{P}}{median}y_u|\]
\end{itemize}
\subsubsection*{Skewness Attributes}
THese are measures of asymmetry in a population. A symmetric distribution of population values should result in a skewness attribute of zero
\begin{itemize}
    \item Pearson's moment coefficient of Skewness: \[a(\mathcal{P}) = \frac{\frac{1}{N}\sum_{u\in\mathcal{P}}(y_u-\overline{y})^3}{[SD_\mathcal{P}(y)]^3}\]    
    \item Pearson's second skewness coefficient (median skewness) given by \[a(\mathcal{P}) = 3\times\frac{(\overline{y}-\underset{u\in\mathcal{P}}{median}y_u)}{SD_\mathcal{P}(y)}\]
    \item Bowley's measure of skewness based on the quartiles \[a(\mathcal{P}) = \frac{(Q_3+Q_1)/2-Q_2}{(Q_3-Q_1)/2}\]
\end{itemize}
\subsubsection*{NAs in R}
Many programs in R accommodate missing data (represented as \textbf{NA}s) 
\begin{itemize}
    \item We can choose to simply omit \textbf{NA}s, we can use the \texttt{na.omit()} function 
\end{itemize}
\subsubsection{Attributes Properties}
\subsubsection*{Location Invariance and Equivariance}
For an attribute $a(\mathcal{P}) = a(y_1,\cdots,y_N)$ we say that for any $m>0$ and $b\in\R$, that the attribute is 
\begin{itemize}
    \item location invariant if \[a(y_1+b,\cdots,y_N+b)=a(y_1,\cdots,y_N)\]
    \item location equivariant if \[a(my_1+b,\cdots,my_N+b)=ma(y_1,\cdots,y_N)+b\]
\end{itemize} 
Population average is location equivariant 
\subsubsection*{Scale Invariance and Equivariance}
For an attribute $a(\mathcal{P}) = a(y_1,\cdots,y_N)$ we say that for any $m>0$ and $b\in\R$, that the attribute is
\begin{itemize}
    \item scale invariant if \[a(m\times y_1,\cdots,m\times y_N)=a(y_1,\cdots,y_N)\]
    \item scale equivariant if \[a(m\times y_1,\cdots,m\times y_N)=m\times a(y_1,\cdots,y_N)\]
    \item location-scale invariant if it is both location invariant and scale invariant \[a(m\times y_1+b,\cdots,m\times y_N+b)=a(y_1,\cdots,y_N)\]
    \item location-scale equivariant if it is both location equivariant and scale equivariant \[a(m\times y_1+b,\cdots,m\times y_N+b)=m\times a(y_1,\cdots,y_N)+b\]
\end{itemize}
Population variance is location-scale equivariant
\subsubsection*{Replication}
The attribute $a(\mathcal{P})$ is 
\begin{itemize}
    \item replication invariant whenever \[a(\mathcal{P}^k)=a(\mathcal{P})\] 
    \item replication equivariant whenever \[a(\mathcal{P}^k)=k\times a(\mathcal{P})\]
\end{itemize}
Population average is replication invariant 
\subsubsection{Influence, Sensitivity Curves, and Breakdown Points}
\subsubsection*{Influence}
If we remove variate $y_u$ then the influence of that variate on the population attribute is quantified by 
\[\tri(a,u)=a(y_1,\cdots,y_{u-1},y_u,y_{u+1},\cdots,y_N)-a(y_1,\cdots,y_{u-1},y_{u+1},\cdots,y_N)\]
Ideally, no signle unit's value should have greater influence than others
\subsubsection*{Sensitivity Curves}   
We can examine the effect on an attribute when we add a variate. To examine the effect,
\begin{itemize}
    \item suppose we have a population of size $N-1$ 
    \item add a variate with the value $y$
    \item new population with $N$ elements is $\{y_1,\cdots,y_{N-1},y\}$
\end{itemize}
Define the sensitivity curve of an attribute as 
\[SC(y;a(\mathcal{P})) = N[a(y_1,\cdots,y_{N-1},y)-a(y_1,\cdots,y_{N-1})]\]
General sensitivity curve function in R can be defined as 
\begin{verbatim}
    sc = function(y.pop, y, attr, ...) {
        N <- length(y.pop) + 1
        sapply( y, function(y.new) {
            N * attr(c(y.new, y.pop), ...) - N * attr(y.pop, ...)
        })
    }
\end{verbatim}
\subsubsection*{Breakdown Points}
Gives an assessment of jusst how large a proportion of the data must be contaminated before the statistic breaks down (becomes useless) \\
The breakdown point of a statistic is the smallest possible fraction of the observations that can be changed to somethings very extreme to make the error large (infinite) \\
Attributes with high breakdown points are called resistant or robust 
\subsubsection{Graphical Attributes}
\subsubsection*{Histograms}
Helps determine how the values are concentrated \\
We can define bins in two ways 
\begin{itemize}
    \item bins of equal size 
    \item bins with equal number of elements but varying size 
\end{itemize}
\subsubsection*{Rules for Number of Bins}
\begin{itemize}
    \item Sturges' rule: \[\text{Number of bins} = \lceil \log_2(N) + 1\rceil\]
    \item Freedman-Diaconis rule: \[\text{Bin size} = 2\times\frac{IQR}{N^{1/3}}\]
    \item Scott's rule: \[\text{Bin size} = 3.5\times\frac{SD}{N^{1/3}}\]
\end{itemize}
\subsubsection*{Scatter-Plots}
Used to see whether two variates $x$ and $y$ are related in some way 
\subsubsection{Power Transformations}
For variates $y$, sometimes can be helpful to re-express the values in a non-linear way (with transformation $T(y)$) so that they are easier to define, understand, or simply to determine \\
A commonly used transformation when $y>0$ is power transformation, which is indexed by a power $\alpha$. The general form is 
\begin{align*}
    T_\alpha(y) &=
    \begin{cases}
        y^\alpha &\alpha\neq 0 \\
        \log(y) &\alpha=0
    \end{cases}
\end{align*}
There transformations are monotonic 
\[y_u<y_v\iff T_\alpha(y_u)<T_\alpha(y_v)\]
What changes, often dramatically, is the relative positions of the variate values \\
The most common purpose of a transformation is to change the shape of the histogram so that it is more symmetric 
Another power transformation is 
\begin{align*}
    T_\alpha(y) &=
    \begin{cases}
        y^\alpha &\alpha>0 \\
        \log(y) &\alpha=0 \\
        -(y)^\alpha &\alpha<0
    \end{cases}
\end{align*}
\subsubsection*{Bump Rule 1: Making Histograms More Symmetric}
\begin{itemize}
    \item If the bump is on "lower" values, then move the power "lower" on the ladder 
    \item If the bump is on "higher" values, then move the power "higher" on the ladder 
\end{itemize}
\subsubsection*{Bump Rule 2: Straightening Scatter-Plots}
Since use of power transformations on both $x$ and $y$, we will have two different ladders of transformations 
\begin{itemize}
    \item $x$ ladder 
    \item $y$ ladder 
\end{itemize}
\begin{figure}[tbhp]
	\begin{center}
		\includegraphics[width=0.6\textwidth]{scatter plot.png}
	\end{center}
\end{figure}
\newpage 
\subsubsection{Order, Rank, and Quantiles}
Recall the order statistics:
\[y_{(1)}\leq y_{(2)}\leq \cdots \leq y_{(N)}\]
which are the ordered values (including ties) of the variate values $y_u\in\mathcal{P}$ \\
The rank statistics (index of $y_i$)
\[r_1,r_2,\cdots,r_N\]
which are the ranks of the variate values $y_1,y_2,\cdots,y_N$ from the $y_u\in\mathcal{P}$ 
\subsubsection*{Quantiles}
\begin{itemize}
    \item $p\in\{\frac{1}{N}, \frac{2}{N}, \cdots, 1\}$ and 
    \item $Q_y(p)$ is the $p$-th quantile of $y$ \[Q_y(p)=y_{(N\times p)}\] is called the quantile function of $y$ for all $p\in[\frac{1}{N}, 1]$
    \item the quantile $Q_y(p)$ for any $p$ locates the variate valuues in the population, and is thus a measure of location 
    \item most (but not all) location measures try to capture central tendency 
\end{itemize}
\subsubsection*{Quantiles that Measure Center}
\begin{itemize}
    \item the median: $Q_y(1/2)$ 
    \item the mid-hinge (average of the first and third quantiles): $\dfrac{Q_y(1/4)+Q_y(3/4)}{2}$
    \item the mid range (average of the minimum and maximum): $\dfrac{Q_y(1/N)+Q_y(1)}{2}$
    \item the trimean: $\dfrac{Q_y(1/4)+2\times Q_y(1/2)+Q_y(3/4)}{4}$
\end{itemize}
\subsubsection*{Quantiles that Measure Spread}
The quantile function can be used to provide some natural measures of spread for the variate $y$:
\begin{itemize}
    \item the range: $Q_y(1)-Q_y(1/N)$
    \item the inter-quantile rnage: $IQR_y=Q_y(3/4)-Q_y(1/4)$
    \item the central $100\times p\%$ range
    \item the difference between any two quantiles might be divided by the difference in the corresponding $p$ values, which is the slope 
\end{itemize}
\subsubsection*{Concentration in Quantiles Plots}
We draw boxes of fixed height to see how many elements are within the box $\Rightarrow$ how concentration changes with $p$
\subsection{Implicit Attributes}
\subsubsection{The Minimum of a Function}
We want the value $\hat{\theta}$ which satisfies 
\[\hat{\theta} = \underset{\theta}{argmin} \quad \rho(\theta;\mathcal{P})\]
where the possible values of $\theta$ may be constrained to be in some set $\Theta$
The maximizing function is the same as minimizing the negation, so we only need to consider minimization \\
The most common form for $\rho(\theta,\mathcal{P})$ is a sum of functions $\rho(\theta,u)$ evaluated at each unit $u\in\mathcal{P}$
\[\rho(\theta,\mathcal{P}) = \sum_{u\in\mathcal{P}}\rho(\theta,u)\]
\subsubsection{Dealing with Influential Units in Linear Regression}
By applying power transformations to the variates, we can see variates that do not behave as the others do. It is likely 
that these will have a large influence on the fitted regression line \\
Thus we can either
\begin{itemize}
    \item remove the units
    \item assign weights to the observations according to their variation 
    \item use a method to find the regression line which is robust to potential outliers
\end{itemize}
\subsubsection*{Weighted Linear Regression}
In weighted linear regression (WLS), the fitted line minimizes the following objective function 
\[(\hat{\alpha},\hat{\beta}) = \underset{\alpha,\beta}{argmin} \quad \sum_{u\in\mathcal{P}}w_u(y_u-\alpha-\beta (x_u - c))^2\]
We assume the weight $w_u$ are known \\
It is common to either set $c=0$ or define $c$ to be the weighted average of the $x_u$ values 
\[c=\overline{x}_w = \frac{\sum_{u\in\mathcal{P}}w_u x_u}{\sum_{u\in\mathcal{P}}w_u}\]
We estimate $\alpha$ and $\beta$ as 
\[\hat{\alpha} = \overline{y}_w - \hat{\beta}(\overline{x}_w - c)\]
\[\hat{\beta} = \frac{\sum_{u\in\mathcal{P}}w_u (x_u - \overline{x}_w)(y_u - \overline{y}_w)}{\sum_{u\in\mathcal{P}}w_u (x_u - \overline{x}_w)^2}\]
where $\overline{y}_w$ and $\overline{x}_w$ are the weighted averages of $y_u$ and $x_u$ respectively
\subsubsection*{Robust Regression}
Points that are far from the linear line should have less weight in the objective function. The objective function will be 
\[(\hat{\alpha},\hat{\beta}) = \underset{\alpha,\beta}{argmin} \quad \sum_{u\in\mathcal{P}}\rho(y_u-\alpha-\beta (x_u - c))\]
In robust regression we modify the loss function $\rho(y_u-\alpha-\beta (x_u - c))$ so that 
\begin{itemize}
    \item gives lower weight than least squares to units with large residuals 
    \item it is quadratice near $0$ and hence behaves similarly to LS for units with small residuals 
\end{itemize}
\subsubsection*{Huber Loss Function}
It achieves these goals by combining the quadratice and absolute value functions 
\begin{align*}
    \rho_k(r) &= 
    \begin{cases}
        \frac{1}{2}r^2 & |r|<k \\
        k(|r|-k/2) & |r|\geq k
    \end{cases}
\end{align*}
Another form of robust regressino involves defining the loss function in terms of least absolute deviations (LAD)
\[(\hat{\alpha},\hat{\beta}) = \underset{\alpha,\beta}{argmin} \quad \sum_{u\in\mathcal{P}}|y_u-\alpha-\beta (x_u - c)|\]
\subsubsection{Gradient Descent}
\subsubsection*{Direction and Step Size}
If $\rho(\theta;\mathcal{P})$ is a differentiable function of $\theta = (\theta_1,\cdots,\theta_p)\in\R^k$, then we can calculate the gradient for any value of $\theta$
\[g = g(\theta) = \tri\rho(\theta;\mathcal{P}) = \cv{\dfrac{\partial \rho(\theta;\mathcal{P})}{\partial \theta_1}\\\vdots\\\dfrac{\partial \rho(\theta;\mathcal{P})}{\partial \theta_k}}\]
By definition, the normalized gradient \[d_i = \dfrac{g_i}{||g_i||}\]    
provides the direction in which $\rho(\theta;\mathcal{P})$ increases or decreases fastest 
\begin{itemize}
    \item $d_i$ indicates the direction of steepest ascent 
    \item $-d_i$ incidates the directino of steepest descent 
\end{itemize}
We iterate and obtain a new estimate of $\theta$ by 
\begin{itemize}
    \item moving in the direction of $-d_i$ and 
    \item taking a step of size $\lambda_i > 0$ \[\hat{\theta}_{i+1} = \hat{\theta}_i - \lambda_id_i\]    
\end{itemize}
Note that the step size $\lambda$ at each iteration can been chosen in a variety of ways 
\begin{enumerate}
    \item could choose a fied value for all $i$ such as $\lambda_i = 0.1$
    \item could define a fixed sequence such as $\lambda_i=0.1 + 1/i$
    \item could perform a line search and algorithmically choose the value of $\lambda_i$ that minimizes \[\rho(\hat{\theta}_i - \lambda_id_i)\] 
\end{enumerate}
\subsubsection*{The Gradient Descent Algorithm}
Given some initial values $\hat{\theta}_0$ 
\begin{enumerate}
    \item initialize $i=0$
    \item Loop until convergence:
    \begin{enumerate}[a.]
        \item Gradient: \[g_i = \nabla\rho(\hat{\theta}_i;\mathcal{P})\]
        \item Gradient direction: \[d_i = \frac{g_i}{||g_i||}\]
        \item Line search: Find the step size $\hat{\lambda}_i$ \[\hat{\lambda}_i = \underset{\lambda>0}{argmin} \quad \rho(\hat{\theta}_i - \lambda d_i)\]    
        \item Update iterate: \[\hat{\theta}_{i+1} = \hat{\theta}_i - \hat{\lambda}_i d_i\]
        \item Converged?
        \begin{itemize}
            \item if iteratesnot changing, Return 
            \item else $i = i + 1$ and repeat
        \end{itemize}
    \end{enumerate}
    \item Return $\hat{\theta} = \hat{\theta}_i$
\end{enumerate}
\newpage 
\subsubsection*{R Code for Gradient Descent}
\begin{lstlisting}[language=R]
gradientDescent <- function(theta = 0, rhoFn, gradientFn, lineSearchFn, 
    testConvergenceFn, maxIterations = 100, tolerance = 1e-06, relative = FALSE, 
    lambdaStepsize = 0.01, lambdaMax = 0.5) {

    converged <- FALSE
    i <- 0

    while (!converged & i <= maxIterations) {
    g <- gradientFn(theta)      ## gradient
    glength <- sqrt(sum(g^2))   ## gradient direction
    if (glength > 0)
        d <- g / glength

    lambda <- lineSearchFn(theta, rhoFn, d, lambdaStepsize = lambdaStepsize,
                            lambdaMax = lambdaMax)

    thetaNew <- theta - lambda * d
    converged <- testConvergenceFn(thetaNew, theta, tolerance = tolerance,
                                    relative = relative)
    theta <- thetaNew
    i <- i + 1
    }

    ## Return last value and whether converged or not
    list(theta = theta, converged = converged, iteration = i, fnValue = rhoFn(theta))
}
\end{lstlisting}
\subsubsection{Systems of Equations}
\subsubsection*{Scalar Valued Attributes}
The average is the value of $\theta$ that solves 
\[\psi(\theta;\mathcal{P}) = \sum_{u\in\mathcal{P}}\psi(\theta;u)=\sum_{u\in\mathcal{P}}y_u-n\theta=0\]   
The weighted average is the value of $\theta$ which solves 
\[\psi(\theta;\mathcal{P}) = \sum_{u\in\mathcal{P}}\psi(\theta;u)=\sum_{u\in\mathcal{P}}w_u(y_u-\theta)=0\]
The $q$-th quantile is the samllest $\theta$ which solves 
\[\psi(\theta;\mathcal{P})=\sum_{u\in\mathcal{P}}\psi(\theta;u)=\sum_{u\in\mathcal{P}}\frac{1}{N}I(y_u\leq\theta)-q=0\]    
\subsubsection*{Vector Valued Attributes}
Least squares
\[\psi(\theta;\mathcal{P}) = \sum_{u\in\mathcal{P}}\psi(\theta;u) = \sum_{u\in\mathcal{P}}(y_u-\alpha-\beta(x_u-c))\cv{1\\x_u-c}=0\]    
Weighted least squares 
\[\psi(\theta;\mathcal{P}) = \sum_{u\in\mathcal{P}}\psi(\theta;u) = \sum_{u\in\mathcal{P}}w_u(y_u-\alpha-\beta(x_u-c))\cv{1\\x_u-c}=0\]
Robust regression
\[\psi(\theta;\mathcal{P}) = \sum_{u\in\mathcal{P}}\psi(\theta;u) = \sum_{u\in\mathcal{P}}\rho'_k(y_u-\alpha-\beta(x_u-c))\cv{1\\x_u-c}=0\]    
\subsubsection{Newton's Method}
Suppose we have a differentiable function $f(x)$ and we wish to find $x=x^*$ which solves \[f(x)=0\]    
Given initial value $x_0$
\begin{itemize}
    \item use a linear function to approximate $f(x)$ in the vicinity of $x_0$ \[f(x)\approx f(x_0)+f'(x_0)(x-x_0)\]    
    \item then find the root of linear approximation to iterate to the next value of $x$ \[0=f(x_0)+f'(x_0)(x-x_0)\implies x_1=x_0-\frac{f(x_0)}{f'(x_0)}\]
    \item continue until $x_{i+1}$ does not differ much from $x_i$ 
\end{itemize}
For implicitly defined scalar attribute $\theta$, we want to find $\theta\in\Theta$ such that 
\[\psi(\theta;\mathcal{P})=0\]    
\begin{itemize}
    \item Given current guess $\hat{\theta} = \hat{\theta}_i$ a first order approximation is \[\psi(\theta;\mathcal{P})\approx \psi(\hat{\theta}_i;\mathcal{P})+\psi'(\hat{\theta}_i;\mathcal{P})\times(\theta-\hat{\theta}_i)\]    
    \item Then the update is the root of the linear approximation and is given by \[\hat{\theta}_{i+1}=\hat{\theta}_i-\frac{\psi(\hat{\theta}_i;\mathcal{P})}{\psi'(\hat{\theta}_i;\mathcal{P})}\]
\end{itemize}
\subsubsection*{The Newton Algorithm}
Given an initial value $\hat{\theta}_0$
\begin{enumerate}
    \item Initialize: $i=0$
    \item Loop:
    \begin{enumerate}
        \item Update iterate: \[\hat{\theta}_{i+1}=\hat{\theta}_i-\frac{\psi(\hat{\theta}_i;\mathcal{P})}{\psi'(\hat{\theta}_i;\mathcal{P})}\]
        \item Converged?
        \begin{itemize}
            \item if iterates not changing, return 
            \item else $i=i+1$ and repeat
        \end{itemize}
    \end{enumerate}
    \item Return $\hat{\theta}=\hat{\theta}_i$
\end{enumerate}
\subsubsection{The Newton raphson Algorithm}
The multivariate analog of Newton's method  \\ 
Given an initial value $\hat{\theta}_0$ 
\begin{enumerate}
    \item Initialize: $i=0$
    \item Loop:
    \begin{enumerate}
        \item Update iterate: \[\hat{\theta}_{i+1} = \hat{\theta}_i - [\psi'(\hat{\theta}_i;\mathcal{P})]^{-1}\psi(\hat{\theta}_i;\mathcal{P})\]    
        \item Converged?
        \begin{itemize}
            \item if iterates not changing, return
            \item else $i=i+1$ and repeat
        \end{itemize}
    \end{enumerate}
    \item Return $\hat{\theta}=\hat{\theta}_i$
\end{enumerate}
\subsubsection{Iteratively Reweighted Least Squares (IRLS)}
Given initial value $\hat{\theta}_0$
\begin{enumerate}
    \item Initialize: $i=0$
    \item Loop:
    \begin{enumerate}
        \item Construct residuals and weights for all $u\in\mathcal{P}$ \[r_u = y_u - z_u'\theta_i\quad w_u = \frac{\rho'(r_u)}{r_u} \] where $z_u = \cv{1 &x_u-c}'$
        \item Solve the weighted least squares problem \[\sum_{u\in\mathcal{P}}w_u(y_u-z_u'\theta)z_u = 0\]   
        \item Update parameter: \[\hat{\theta}_{i+1} = \hat{\theta}_i\]    
        \item Converged?
        \begin{itemize}
            \item if iterates not changing, return 
            \item else $i=i+1$ and repeat
        \end{itemize}
    \end{enumerate}
    \item return $\hat{\theta}=\hat{\theta}_i$
\end{enumerate}
\newpage 

\section{Samples}
If we have a sample or a subset $S$ of $n<<N$ units, then the attribute $a(S)$ calculated based 
on this sample is an estimate of its population counterpart $a(\mathcal{P})$
\[a(S) = \hat{a(\mathcal{P})} = a(\hat{\mathcal{P}})\]   
\subsubsection*{Sample Error}
Any difference between the actual values of the estimates $a(S)$ and the quantity being estimated (the estimand) $a(\mathcal{P})$ is an error 
\[\text{sample error} = a(S) - a(\mathcal{P})\]
The nature of this error will depend on the sample and attribute 
\subsubsection*{Fisher Consistency}
If the sample $S$ is equal to the population $\mathcal{P}$ then the sample error should be zero (or non-existent), $a(S) = a(\mathcal{P})$ \\
This would mean the estimation is in some sense consistent, called Fisher consistency
\subsection{All Possible Samples}
\subsubsection*{Sample Error}
The sample error for a sample $S$ of size $n$ is 
\[a(S)-a(\mathcal{P}) = \frac{1}{n}\sum_{u\in S}y_u - \frac{1}{N}\sum_{u\in\mathcal{P}}y_u\]
\subsubsection*{Average Sample Error}
The average sample error over all possible samples of size $n$ is 
\[\text{Average Sample Error} = (\frac{1}{M}\sum_{i=1}^{M}a(S_i))-a(\mathcal{P})\]    
\subsubsection{Consistency and the Effect of Sample Size}
The nature of sample error depends largely on the sample size
\begin{itemize}
    \item As sample size increases, the sample approaches the population  
    \item Attribute values will concentrate even more around the population value 
\end{itemize}
We can quantify the concentration by 
\[||a(S)-a(\mathcal{P})|| = ||\frac{1}{n}\sum_{u\in S}y_u-\frac{1}{N}\sum_{u\in\mathcal{P}}y_u||<c\]   
for some $c>0$, we could proportion of samples that satisfy this \\
Consider a population $\mathcal{P}$ of size $N<\infty$
\begin{itemize}
    \item For each $n$, construct the set of all possible samples \[\mathcal{P}_S(n) = \{S:S\subset\mathcal{P}, |S|=n\}\]    
    \item For any $c>0$, \[\mathcal{P}_a(c,n) = \{S:S\subset\mathcal{P}_S(n), ||a(S)-a(\mathcal{P})||<c\}\] and define the proportion \[p_a(c,n) = \frac{|\mathcal{P}_a(c,n)|}{|\mathcal{P}_S(n)|}\]   for all $c>0$ and $n\leq N$
\end{itemize}
\subsubsection{Comparisons Across Attributes}
We use relative absolute sample error to compare different attributes 
\[\mathcal{P}^*_a(c,n)=\{S:S\subset\mathcal{P}_S(n), \frac{||a(S)-a(\mathcal{P})||}{||a(\mathcal{P})||}<c\}\] and define the corresponding proportino, for all $c>0$ and $n\leq N$
\[p_a^*(c,n) = \frac{|\mathcal{P}_a^*(c,n)|}{|\mathcal{P}_S(n)|}\]
\subsection{Selecting Samples}
\subsubsection{Randomly Selecting m Samples}
Consider drawing samples of size $n$ from population $\mathcal{P}$. The population $\mathcal{P}_S$ of all such 
samples has size $M={\ds N\choose n}$ and is denoted \[\mathcal{P}_S=\{S_1,S_2,\cdots,S_M\}\]    
Any attribute $a(S_u)$ is now just a variate on that unit \\
We have a population of attributes 
\[\mathcal{P}_{a(S)} = \{a(S_1),a(S_2),\cdots,a(S_M)\}\]
\subsubsection{Quantifying Sample Error}
In principle we select a sample $S$ from the population $\mathcal{P}_S$ contianing all possible samples 
\begin{itemize}
    \item We do so with some probability $p(S)\geq 0$ of being selected, require \[\sum_{S\in\mathcal{P}_S}p(S)=1\]    
    \item For any sample $S\in\mathcal{P}_S$, we have sample error \[\text{Sample Error}=a(S)-a(\mathcal{P})\]    
    \item Aveerage sample error is \[\text{Average Sample Error} = \frac{1}{M}\sum_{S\in\mathcal{P}_S}[a(S)-a(\mathcal{P})]\]    
\end{itemize}
\subsubsection{Sampling Bias, Sampling Variance, and Sampling Mean Squared Error}
\subsubsection*{Sampling Bias}
\begin{align*}
    \text{Sampling Bias} &= E[a(S)] - a(\mathcal{P}) \\
                         &= \sum_{S\in\mathcal{P}_S}a(S)p(S) - a(\mathcal{P}) \\
                         &= \sum_{S\in\mathcal{P}_S}[a(S)-a(\mathcal{P})]p(S)
\end{align*}
Sampling bias is the expected sample error induced by the repeated random sampling of $S$ from $\mathcal{P}_S$ \\
If sampling bias is zero, then $a(S)$ is unbiased estimator of $a(\mathcal{P})$ 
\subsubsection*{Sampling Variance}
\[Var[a(S)] = E[(a(S)-E(a(S)))^2] \]    
This quantifies dispersion in sample errors
\subsubsection*{Mean Squared Error}
Use to quantify the expected squared distance between $a(S)$ and $a(\mathcal{P})$ 
\[MSE[a(S)] = E[(a(S)-a(\mathcal{P}))^2] = Var[a(S)] + [\text{Sampling Bias}]^2\]    
\subsubsection{Sampling Mechanisms}
A sequence of the first $k$ units $u_i$ selected from $\mathcal{P}$ is 
\[s_k = (u_{i1},u_{i2},\cdots,u_{ik})\]
A sampling mechanism is defined by the probabilities 
\[Pr(u)\text{ and } Pr(u|k,s_{k-1})\]
where 
\begin{itemize}
    \item the first unit is selected with probabilityy $Pr(u)$ and 
    \item the probability of the sequence of the first $k$ units selected is 
    \[Pr(s_k) = Pr(u_{i1})Pr(u_{i2}|2,s_{1})\cdots Pr(u_{ik}|k,s_{k-1})\]
\end{itemize}
\subsubsection*{Simple Random Sampling Without Replacement (SRSWOR)}
The sampling mechanism is \[Pr(u) = \frac{1}{N} \text{ and } Pr(u|k,s_{k-1}) = \frac{1}{N-k+1}\]
The probability of the sequence $s_n$ is 
\[Pr(s_n) = \frac{1}{N}\frac{1}{N-1}\cdots\frac{1}{N-n+1} \]    
which is the same for all $n!$ permutations, so 
\[p(S) = \frac{n!}{N(N-1)\cdots(N-n+1)} = \frac{1}{{N\choose n}}\]
\subsubsection*{Simple Random Sampling With Replacement (SRSWR)}
The sampling mechanism is 
\[Pr(u) = \frac{1}{N} = Pr(u|k,s_{k-1})\]
Sample $S$ can contain one or more replicated units \\
The probability of the sequence $s_n$ is 
\[Pr(s_n) = \frac{1}{N^n}\]
\subsubsection{Unit Inclusion Probabilities}
The inclusion probability for unit $u$ is probability of unit $u$ being included in the sample 
\[\pi_u = Pr(u\in S) = \sum_{S\in\mathcal{P}_S}p(S)\times I(u\in S)\]    
Such inclusion probabilities are of interest to calculate in addition to $p(S)$ and may in fact be derived from $p(S)$ \\
Consider the following random variable 
\begin{align*}
    D_u = \begin{cases}
        1 & \text{if } u\in S \\
        0 & \text{if } u\notin S
    \end{cases}
\end{align*}
The expectation of this random variable is 
\begin{align*}
    E[D_u] &= 1\times Pr(D_u=1) + 0\times Pr(D_u=0) \\
           &= Pr(u\in S) \\
           &= \pi_u
\end{align*}
The variance of this random variable is 
\begin{align*}
    Var[D_u] &= E[D_u^2] - E[D_u]^2 \\
             &= 1^2\times Pr(D_u=1) + 0^2\times Pr(D_u=0) - \pi_u^2 \\
             &= \pi_u - \pi_u^2 \\
             &= \pi_u(1-\pi_u)
\end{align*}
The probability that $u$ and $v$ are in the sample $S$ is called the joint inclusion probability and is given by 
\[\pi_{uv} = Pr(u\in S, v\in S)\]
The relationship is 
\begin{align*}
    Cov(D_u,D_v) &= E[D_u\times D_v] - E[D_u]E[D_v] \\
               &= Pr(u\in S, v\in S) - \pi_u\pi_v \\
               &= \pi_{uv} - \pi_u\pi_v
\end{align*}
\subsubsection*{Simple Random Sampling Without Raplacement (SRSWOR)}
The inclusion probability is 
\[\pi_u = Pr(u\in S) = \frac{n}{N}\]
The joint inclusion probability is
\[\pi_{uv} = Pr(u\in S, v\in S) = \frac{n(n-1)}{N(N-1)}\]
\subsubsection*{Simple Random Sampling Withou Replacement (SRSWR)}
The inclusion probability is 
\[\pi_u = 1 - \left(\frac{N-1}{N}\right)^n\]
The joint probability is 
\[\pi_{uv} = 1 - 2\left(\frac{N-1}{N}\right)^n + \left(\frac{N-2}{N}\right)^{n}\]
\subsubsection{Estimating Tools}
A natural estimate of a population total 
\[a(\mathcal{P}) = \sum_{u\in\mathcal{P}}y_u\]
is the Horvitz-Thompson estimate defined as 
\[\hat{a}(\mathcal{P}) = a_{HT}(S) = \sum_{u\in S}\frac{y_u}{\pi_u}\]
where the contribution for each unit in the sample is weighted inversely by $\pi_u$, its probability of inclusion in $S$
\begin{itemize}
    \item if the probability of inclusion is small, then the weight will be high
    \item if the probability of inclusion is large, then the weight will be low 
\end{itemize}
\subsubsection*{The Horvitz-Thompson Estimator}
\subsubsection*{Bias, Variance, and Mean Squared Error}
The bias of the HT estimator 
\begin{align*}
    E[\tilde{a}_{HT}(S) - a(\mathcal{P})] &= 0
\end{align*}
The variance of the HT estimator is 
\[Var[\tilde{a}_{HT}(S)] = \sum_{u\in P}\sum_{v\in P}(\pi_{uv} - \pi_u\pi_v)\frac{y_u}{\pi_u}\frac{y_v}{\pi_v} = \sum_{u\in S}\sum_{v\in S}Cov(D_u,D_v)\frac{y_u}{\pi_u}\frac{y_v}{\pi_v}\] 
The HT estimator is unbiased, the mean squared error is equal to the variance 
\[MSE = Var\]    
The HT estimate of the variance of the HT estimator is 
\[\hat{Var}[\tilde{a}_{HT}(S)] = \sum_{u\in S}\sum_{v\in S}(\frac{\pi_{uv}-\pi_u\pi_v}{\pi_{uv}})\frac{y_u}{\pi_u}\frac{y_v}{\pi_v}\]    
Using HT estimation we are able to construct 
\begin{itemize}
    \item an estimate of population total 
    \item an estimate of the variance of this estimator 
    \item both estimators are unbiased
\end{itemize}
\subsubsection{Sampling Designs}
The pair $(\mathcal{P}_S,p(S))$ is called a sampling design 
\begin{itemize}
    \item Together they determine which samples are possible and with what probability they are selected 
    \item The SRSWOR, SRSWR and SRSWH frameworks provide examples of different sampling designs 
    \item The sampling design is ours to choose 
    \begin{itemize}
        \item We may choose $\mathcal{P}_S$ so that the values $a(S)$ for $S\in\mathcal{P}_S$ are constrained to be near $a(\mathcal{P})$ 
        \item We may choose $p(S)$ so that samples $S\in\mathcal{P}_S$ that have $a(S)$ close to $a(\mathcal{P})$ have higher probability $p(S)$ of being selected 
    \end{itemize}
\end{itemize}

\newpage 
\section{Inference}
\subsection{Sources of Error}
\subsubsection{Target and Study Populations}
The target population is the population about which we would like to draw an inference \\
The study population is the population from which samples are taken \\
The difference between attribute evaluated on the two populations is the \textbf{study error}
\[Study\ Error = a(\mathcal{P}_{study}) - a(\mathcal{P}_{target})\]
Then if we obtain a sample $S$ to draw inferences about the target popiulation $\mathcal{P}_{target}$, the error for a 
given attribute $a(\cdot)$ is 
\begin{align*}
    a(S) - a(\mathcal{P}_{target}) &= [a(S) - a(\mathcal{P}_{study})] + [a(\mathcal{P}_{study}) - a(\mathcal{P}_{target})] \\
    &= Sample\ Error + Study\ Error
\end{align*}
\subsubsection*{Future Realizations}
The units which are not available at the time of study 
\subsubsection{Measurement Error}
Every measuring system has at least three sources of potential error:
\begin{itemize}
    \item the measuring device (gauge)
    \item the person reading or recording the measurement (operator)
    \item the method followed to take the measurement 
\end{itemize}

\subsection{Comparing Sub-Populations}
\subsubsection{Anatomy of a Significance Test}
We would like to quantify, numerically, how unusual the difference between $a(\mathcal{P}_1)$ and $a(\mathcal{P}_2)$ is relative to randomly mixed sub-populations \\
Following steps are used to gather such evidence:
\begin{enumerate}
    \item Suppose the sub-populations were randomly drawn from the same population, known as the \textbf{null hypothesis}
    \item Construct a \textbf{discrepency measure} that quantifies how inconsistent our data is with the null hypothesis, large value indicate evidence against the null hypothesis 
    \item Obtain the \textbf{observed discrepency} by calculating the discrepency measure on the observed / unshuffled sub-populations 
    \item Obtain the \textbf{observed p-value} by calculating the probability that a randomly shuffled sub-population has a discrepency measure at least as large as the observed discrepency, small values indicate evidence against the null hypothesis 
\end{enumerate}
\subsubsection*{Null Hypothesis}
Each of the following (equivalent) statements constitutes the null hypothesis we are testing:
\begin{itemize}
    \item $H_0$: the sub-populations $\mathcal{P}_1$ and $\mathcal{P}_2$ were randomly drawn from the same population 
    \item $H_0$: $\mathcal{P}_1$ and $\mathcal{P}_2$ were created by randomly assigning units in the same population to one of the two sub-populations 
    \item $H_0$: $\mathcal{P}_1$ and $\mathcal{P}_2$ were generated by random mixing 
\end{itemize}
\subsubsection*{Discrepency Measure}
A discrepency measure $D(\mathcal{P}_1, \mathcal{P}_2)$ quantifies how inconsistent data is with the null hypothesis, is defined so that large values indicate evidence against the null hypothesis \\
Form of $D(\mathcal{P}_1, \mathcal{P}_2)$ depends on how we want to compare $\mathcal{P}_1$ and $\mathcal{P}_2$ 
\begin{itemize}
    \item compare measures of location, typically based on difference \[a(\mathcal{P}_1) - a(\mathcal{P}_2)\]
    \item compare measures of spread, typically based on ratios \[\frac{a(\mathcal{P}_1)}{a(\mathcal{P}_2)}\]
\end{itemize}
\subsubsection*{Observed Discrepency}
The observed discrepency $d_{obs}$, is the value of discrepency measure $D$ calculated on the two observed / unshuffled sub-populations 
\[d_{obs} = D(\mathcal{P}_1, \mathcal{P}_2)\]
\subsubsection*{Observed p-value}
The observed p-value is the probability that a randomly shuffled sub-population has a discrepency measure at least as large as the observed discrepency 
\[p-value = Pr(D\geq d_{obs}|\text{$H_0$ is true})\]
\begin{itemize}
    \item $p-value < 0.001$: very strong evidence against $H_0$
    \item $0.001 < p-value < 0.01$: strong evidence against $H_0$
    \item $0.01 < p-value < 0.05$: moderate evidence against $H_0$
    \item $0.05 < p-value < 0.1$: weak evidence against $H_0$
    \item $p-value > 0.1$: no evidence against $H_0$
    \item in extreme cases, $p-value=0$, the hypothesis must be false 
\end{itemize}
There are $\binom{N_1+N_2}{N_1} = \binom{N_1+N_2}{N_2}$ possible permutations of the observed data \\
Because that's too many permutations to consider, we typically just use $M$ of them, the p-value is approximated as 
\[\frac{1}{M}\sum_{i=1}^{M}I(D(\mathcal{P}_{1,i},\mathcal{P}_{2,i})\geq d_{obs})\]
\subsubsection*{Putting all Together}
Test of Significance Algorithm 
\begin{enumerate}
    \item State the null hypothesis: $H_0$: $\mathcal{P}_1$ and $\mathcal{P}_2$ are drawn from the same population 
    \item Construct a measure of discrepency $D=D(\mathcal{P}_1,\mathcal{P}_2)$ where large values indicate evidence against the null hypothesis 
    \item Calculate the observed discrepency $d_{obs} = D(\mathcal{P}_1,\mathcal{P}_2)$
    \item Shuffle the sub-populations $M$ times and calculate the observed p-value \[p-value = Pr(D\geq d_{obs}|\text{$H_0$ is true})\approx \frac{1}{M}\sum_{i=1}^{M}I(D(\mathcal{P}_{1,i},\mathcal{P}_{2,i})\geq d_{obs})\]
\end{enumerate}
\subsubsection{A t-like Discrepency Measure}
When comparing two sub-populations on the basis of a measure of location, one particularly useful discrepancy measure is 
\[D(\mathcal{P}_1,\mathcal{P}_2) = \frac{a(\mathcal{P}_1)-a(\mathcal{P}_2)}{\sqrt{Var(\tilde{a}(\mathcal{P}_1))+Var(\tilde{a}(\mathcal{P}_2))}}\]
This is scale-invariant 
\subsubsection{Multiple Testing}
\subsubsection*{Estimating $d^*_{obs}$}
\begin{itemize}
    \item Suppose we have $K$ discrepancy measures $D_1,D_2,\cdots,D_K$ 
    \begin{itemize}
        \item The combined discrepancy measure is \[D^* = 1 - p-value_{min}\]
    \end{itemize}
    \item For $i=1,\cdots,M_{inner}$ and each discrepancy $k=1,\cdots,K$
    \begin{itemize}
        \item randomly mix the two sub-populations $\mathcal{P}_1$ and $\mathcal{P}_2$ yielding $\mathcal{P}^*_{1,i}$ and $\mathcal{P}^*_{2,i}$
        \item calculate \[d_{k,i} = D_k(\mathcal{P}^*_{1,i},\mathcal{P}^*_{2,i})\]
        \item we estimate each p-value labelled as $p_k$ with \[\hat{p}_k = \frac{1}{M_{inner}}\sum_{i=1}^{M_{inner}}I(d_{k,i}\geq d_{obs,k})\]
    \end{itemize}
    \item we estimate $p-value_{min}$ as \[\hat{d}^*_{obs} = 1 - \min_{k=1,\cdots,K}\hat{p}_k\]
\end{itemize}
\subsubsection*{Estimating $p-value^*$}
\[\hat{p-value}^* = \frac{1}{M_{outer}}\sum_{i=1}^{M_{outer}}I(D^*(\mathcal{P}^*_{1,i},\mathcal{P}^*_{2,i})\geq d^*_{obs})\]
\subsection{Interval Estimation}
\subsubsection{Random vs. Observed Intervals}
Suppose the attribute of interest is the population average $a(\mathcal{P}) = \mu = \overline{y}$
\[\overline{Y}\sim N(\mu, \frac{\sigma^2}{n}(\frac{N-n}{N-1}))\]
where \[\sigma^2 = \frac{1}{N}\sum_{u\in\mathcal{P}}(y_u - \mu)^2\]
\subsubsection*{Random Intervals}
using the standardized random variable and sepcified $p\in(0,1)$ we can find a constant $c>0$ such that 
\[1-p = Pr(\overline{Y} - c\cdot \frac{\sigma}{\sqrt{n}}\sqrt{\frac{N-n}{N-1}} \leq \mu \leq \overline{Y} + c\cdot \frac{\sigma}{\sqrt{n}}\sqrt{\frac{N-n}{N-1}})\]
\subsubsection*{Observed Intervals}
\[\left[\overline{y} - c\cdot \frac{\sigma}{\sqrt{n}}\sqrt{\frac{N-n}{N-1}} \leq \mu \leq \overline{y} + c\cdot \frac{\sigma}{\sqrt{n}}\sqrt{\frac{N-n}{N-1}}\right]\]
where $\overline{y}$ is the sample mean and $\sigma$ is the sample standard deviation
\subsubsection{Student $t$ Based Intervals}
\subsubsection*{Standard Error vs. Standard Deviation}
The standard error is an estimate of the standard deviation of the corresponding estimator 
\[SE[\tilde{a}(\mathcal{P})] = \hat{SD}[\tilde{a}(\mathcal{P})]\]
\[\frac{\overline{Y}-\mu}{\frac{\tilde{\sigma}}{\sqrt{n}}\sqrt{\frac{N-n}{N-1}}}\sim t_{n-1}\]
In the special case that $a(\mathcal{P}) = \overline{Y}$ is the population average then the standard deviation is 
\[SD[\tilde{a}(S)] = \frac{\sigma}{\sqrt{n}}\sqrt{\frac{N-n}{N-1}}\]
the standard deviation estimator is 
\[\tilde{SD}[\tilde{a}(S)] = \frac{\tilde{\sigma}}{\sqrt{n}}\sqrt{\frac{N-n}{N-1}}\]
the standard error is 
\[\hat{SE}[\tilde{a}(S)] = \frac{\hat{\sigma}}{\sqrt{n}}\sqrt{\frac{N-n}{N-1}}\]

\subsection{Resampling}
\subsubsection{Resampling Intuitions}
In practice, the population from which our sample was taken cannot generally be repeatedly sampled, we only have one sample \\
\textbf{Resampling} methods aim to mimic this process by repeatedly sampling $S$ as if it were $\mathcal{P}$ \\
\textbf{Bootstrapping} is a resampling method that sample with replacement from the sample $S$ to create a new sample $S^*$ of the same size as $S$ 
\subsubsection{Bootstrap}
The bootstrap esetimate is, at best, as good as $a(S)$ 
\[bootstrap\ sample\ error = a(S^*) - a(\mathcal{S})\]
\subsubsection*{Bootstrap Standard Deviation}
The standard deviation of the corresponding sample attribute's estimator can be estimated from the bootstrap distribution with 
\[\hat{SD}_*[\tilde{a}(S^*)] = \sqrt{\frac{\sum_{b=1}^{B}(a(S^*_b)-\overline{a}^*)^2}{B-1}}\]
\subsubsection{Bootstrap Confidence Intervals}
A $95\%$ naive normal theory bootstrap confidence interval for a population attribute $a(\mathcal{P})$ is 
\[a(S)\pm 1.96\cdot \hat{SD}_*[\tilde{a}(S^*)]\]
\subsubsection{Bootstrap t Confidence Intervals}



\newpage 
\section{Prediction}
\subsection{Accuracy of Prediction}
\subsubsection*{Prediction}
The response model \[y = \mu(x) + error\]
yields the predictor function $\hat{\mu}(x)$
\subsubsection*{Accuracy of Prediction}
One measure of inaccuracy over the population $\mathcal{P}$ (of size $N$) is the average prediction squared error (APSE)
\[\frac{1}{N}\sum_{u\in\mathcal{P}}(y_u - \hat{\mu}(x_u))^2\]



\subsection{Predictions over Multiple Samples}
\subsubsection{Calculating APSE Over Many Samples}
The inaccuracy of a predictor function $\hat{\mu}(x)$ is measured by its average prediction squared error 
\[APSE(\mathcal{P},\hat{\mu}_S) = \frac{1}{N}\sum_{u\in\mathcal{P}}(y_u - \hat{\mu}_S(x_u))^2\]
The average APSE over all $M$ samples should be a better measure of the quality of a predictor functino 
\[APSE(\mathcal{P},\tilde{\mu}) = \frac{1}{M}\sum_{i=1}^{M}APSE(\mathcal{P},\hat{\mu}_{S_i}) = \frac{1}{M}\sum_{i=1}^M \frac{1}{N}\sum_{u\in\mathcal{P}}(y_u - \hat{\mu}_{S_i}(x_u))^2\]
\subsubsection{Decompose $APSE(\mathcal{P}, \tilde{\mu})$}
\begin{align*}
    APSE(\mathcal{P},\tilde{\mu}) &= \frac{1}{N}\sum_{u\in\mathcal{P}}(y_u - \tau(x_u))^2 + \frac{1}{M}\sum_{j=1}^{M}\frac{1}{N}\sum_{u\in\mathcal{P}}(\hat{\mu}_{S_j}(x_u) - \overline{\hat{\mu}}(x_u))^2 + \frac{1}{N}\sum_{u\in\mathcal{P}}(\overline{\hat{\mu}}(x_u) - \tau(x_u))^2 \\
    &= Ave_X(Var[Y|X]) + Var[\tilde{\mu}] + Bias^2[\tilde{\mu}] 
\end{align*}




\subsection{Back to Reality}


\end{document}
